% titel: Blatt 8, A2: Newton-Verfahren
% skript: Kap. 4.3 (Newton-Verfahren, Satz 4.3.8)

% Teil I: Werkzeuge, kein Lösungsweg.

Das Newton-Verfahren von Hand und per Code anwenden, die
Konvergenz-Voraussetzungen (Satz 4.3.8) prüfen und eine Nullstelle mit hoher
Genauigkeit berechnen.

\begin{werkzeugkasten}[Newton-Iteration]
Für eine differenzierbare Funktion $f$ mit Nullstelle $x^*$: Start $x_0$
nahe der Nullstelle, dann wiederholt
\begin{equation*}
  x_{k+1}=x_k-\frac{f(x_k)}{f'(x_k)}.
\end{equation*}
Die Linearisierung von $f$ bei $x_k$ wird auf $0$ gelesen. Unter den
Voraussetzungen von Satz 4.3.8 (Start im ``Anziehungskegel'', passende
Monotonie/Konvexität) konvergiert die Folge quadratisch gegen $x^*$.

\textit{Skript:} Kap. 4.3 (Satz 4.3.8).
\end{werkzeugkasten}

\begin{werkzeugkasten}[Voraussetzungen von Satz 4.3.8 prüfen]
Für $f(x)=x^5-2x-1$ auf $[1,2]$:
\begin{itemize}
  \item $f(1)=-2<0$, $f(2)=27>0$: Vorzeichenwechsel, Nullstelle vorhanden.
  \item $f'(x)=5x^4-2>0$ auf $[1,2]$ (strict wachsend): Nullstelle
        \emph{eindeutig}.
  \item $f''(x)=20x^3>0$ auf $[1,2]$: konvex; für Start $x_0=2$ (auf der
        ``steigenden'' Seite) greifen die Bedingungen.
\end{itemize}
\textit{Skript:} Kap. 4.3, Satz 4.3.8.
\end{werkzeugkasten}

\begin{achtung}[Darauf achten]
\begin{itemize}
  \item Beim Handschritt die Ableitung korrekt bilden
        ($f'=5x^4-2$), nicht nur $f$ auswerten.
  \item $f'(x)\ne0$ in der Nähe der Nullstelle prüfen, sonst teilt man
        durch Null.
  \item Bei der Ex-Gleichung $e^x=2-x^2$ als $g(x)=e^x+x^2-2$ umformen und
        auf $[0,1]$ starten.
  \item Beim Code: die Iteration in eine Schleife, Abbruch nach genügend
        vielen Schritten (Genauigkeit steigt quadratisch); Werte stimmen,
        wenn $|f(x_k)|$ klein wird.
\end{itemize}
\end{achtung}
